\documentclass[12pt,a4paper]{article}

\usepackage[T2A]{fontenc}
\usepackage[utf8]{inputenc}
\usepackage[russian]{babel}
\usepackage{amsmath,amssymb,amsthm}
\usepackage{enumitem}
\usepackage{listings}
\usepackage{xcolor}
\usepackage{tikz}
\usepackage{hyperref}
\usepackage{geometry}

\geometry{margin=2cm}

\lstdefinestyle{code}{
    basicstyle=\ttfamily\small,
    backgroundcolor=\color{gray!10},
    frame=single,
    breaklines=true,
    columns=fullflexible,
    keepspaces=true,
    showstringspaces=false,
    literate=
      {А}{{\CYRA}}1
      {Б}{{\CYRB}}1
      {В}{{\CYRV}}1
      {Г}{{\CYRG}}1
      {Д}{{\CYRD}}1
      {Е}{{\CYRE}}1
      {Ё}{{\CYRYO}}1
      {Ж}{{\CYRZH}}1
      {З}{{\CYRZ}}1
      {И}{{\CYRI}}1
      {Й}{{\CYRISHRT}}1
      {К}{{\CYRK}}1
      {Л}{{\CYRL}}1
      {М}{{\CYRM}}1
      {Н}{{\CYRN}}1
      {О}{{\CYRO}}1
      {П}{{\CYRP}}1
      {Р}{{\CYRR}}1
      {С}{{\CYRS}}1
      {Т}{{\CYRT}}1
      {У}{{\CYRU}}1
      {Ф}{{\CYRF}}1
      {Х}{{\CYRH}}1
      {Ц}{{\CYRC}}1
      {Ч}{{\CYRCH}}1
      {Ш}{{\CYRSH}}1
      {Щ}{{\CYRSHCH}}1
      {Ъ}{{\CYRHRDSN}}1
      {Ы}{{\CYRERY}}1
      {Ь}{{\CYRSFTSN}}1
      {Э}{{\CYREREV}}1
      {Ю}{{\CYRYU}}1
      {Я}{{\CYRYA}}1
      {а}{{\cyra}}1
      {б}{{\cyrb}}1
      {в}{{\cyrv}}1
      {г}{{\cyrg}}1
      {д}{{\cyrd}}1
      {е}{{\cyre}}1
      {ё}{{\cyryo}}1
      {ж}{{\cyrzh}}1
      {з}{{\cyrz}}1
      {и}{{\cyri}}1
      {й}{{\cyrishrt}}1
      {к}{{\cyrk}}1
      {л}{{\cyrl}}1
      {м}{{\cyrm}}1
      {н}{{\cyrn}}1
      {о}{{\cyro}}1
      {п}{{\cyrp}}1
      {р}{{\cyrr}}1
      {с}{{\cyrs}}1
      {т}{{\cyrt}}1
      {у}{{\cyru}}1
      {ф}{{\cyrf}}1
      {х}{{\cyrh}}1
      {ц}{{\cyrc}}1
      {ч}{{\cyrch}}1
      {ш}{{\cyrsh}}1
      {щ}{{\cyrshch}}1
      {ъ}{{\cyrhrdsn}}1
      {ы}{{\cyrery}}1
      {ь}{{\cyrsftsn}}1
      {э}{{\cyrerev}}1
      {ю}{{\cyryu}}1
      {я}{{\cyrya}}1
}

\lstset{style=code}

\title{Билет №56 \\ \small Конфигурационное управление. Системы управления версиями. Централизованное и распределённое управление версиями. Git Workflow, GitHub Workflow. Системы CI/CD. (СпБГУ)}
\author{Сериков Владислав Александрович}
\date{16 мая 2026}

\begin{document}

\maketitle
\tableofcontents
\newpage

\section{Конфигурационное управление}

\subsection{Общее определение}

\textbf{Конфигурационное управление} --- это дисциплина управления составом, состоянием и изменениями сложной системы на протяжении её жизненного цикла.

В разработке программного обеспечения под конфигурационным управлением обычно понимают совокупность процессов, инструментов и правил, которые позволяют:

\begin{enumerate}
    \item однозначно определить, из каких компонентов состоит программная система;
    \item фиксировать состояния системы в контрольных точках;
    \item управлять изменениями;
    \item обеспечивать воспроизводимость сборок и поставок;
    \item отслеживать авторство, причины и историю изменений;
    \item восстанавливать предыдущие версии системы;
    \item контролировать согласованность кода, документации, тестов, конфигураций и окружений.
\end{enumerate}

\subsection{Конфигурационная единица}

\textbf{Конфигурационная единица} --- это элемент системы, который помещается под управление конфигурацией и изменения которого необходимо контролировать.

Примеры конфигурационных единиц:

\begin{itemize}
    \item исходный код;
    \item исполняемые файлы;
    \item библиотеки;
    \item файлы конфигурации;
    \item схемы баз данных;
    \item тестовые данные;
    \item техническая документация;
    \item Dockerfile, Helm-чарты, Terraform-манифесты;
    \item спецификации API;
    \item скрипты сборки и развёртывания.
\end{itemize}

\subsection{Версия, ревизия, релиз}

\textbf{Версия} --- зафиксированное состояние некоторого объекта в определённый момент времени.

\textbf{Ревизия} --- конкретное изменение или конкретное состояние объекта в истории системы управления версиями.

\textbf{Релиз} --- версия программной системы, предназначенная для передачи пользователю, заказчику или эксплуатации.

Например:

\begin{lstlisting}
v1.0.0  -- первый стабильный релиз;
v1.1.0  -- добавлена новая функциональность;
v1.1.1  -- исправлены ошибки без изменения функциональности.
\end{lstlisting}

Часто используется схема семантического версионирования:

\[
MAJOR.MINOR.PATCH
\]

где:

\begin{itemize}
    \item \texttt{MAJOR} увеличивается при несовместимых изменениях API;
    \item \texttt{MINOR} увеличивается при добавлении новой совместимой функциональности;
    \item \texttt{PATCH} увеличивается при исправлении ошибок.
\end{itemize}

\subsection{Базовая линия}

\textbf{Базовая линия} --- формально зафиксированное и согласованное состояние системы или её части, которое служит основанием для дальнейшей разработки.

Примеры базовых линий:

\begin{itemize}
    \item версия требований, утверждённая заказчиком;
    \item стабильная версия архитектурной документации;
    \item состояние исходного кода перед началом тестирования;
    \item релиз-кандидат;
    \item промышленный релиз.
\end{itemize}

Базовая линия важна тем, что после её утверждения изменения должны проходить контролируемую процедуру согласования.

\subsection{Управление изменениями}

\textbf{Управление изменениями} --- процесс регистрации, анализа, согласования, реализации и проверки изменений конфигурационных единиц.

Типичный процесс управления изменением:

\begin{enumerate}
    \item Возникает потребность в изменении.
    \item Создаётся запрос на изменение.
    \item Запрос анализируется.
    \item Оцениваются риски, стоимость и влияние на систему.
    \item Принимается решение: принять, отклонить или отложить изменение.
    \item Изменение реализуется.
    \item Изменение тестируется.
    \item Изменение включается в базовую линию или релиз.
\end{enumerate}

\subsection{Задачи конфигурационного управления}

Основные задачи конфигурационного управления:

\begin{enumerate}
    \item \textbf{Идентификация конфигурации.}

    Определение состава системы и конфигурационных единиц.

    \item \textbf{Контроль изменений.}

    Управление тем, какие изменения могут быть внесены, кем и на каком основании.

    \item \textbf{Учёт состояния конфигурации.}

    Хранение информации о версиях, изменениях, релизах, дефектах, запросах.

    \item \textbf{Аудит конфигурации.}

    Проверка соответствия фактического состояния системы заявленному.

    \item \textbf{Управление сборками и релизами.}

    Обеспечение воспроизводимого процесса сборки, тестирования и поставки.
\end{enumerate}

\subsection{Связь конфигурационного управления с DevOps}

В современных процессах разработки конфигурационное управление тесно связано с DevOps-практиками:

\begin{itemize}
    \item инфраструктура описывается как код;
    \item окружения создаются автоматически;
    \item конфигурации хранятся в системах управления версиями;
    \item сборка и поставка выполняются через CI/CD;
    \item релизы фиксируются тегами;
    \item состояние системы можно восстановить из репозитория.
\end{itemize}

\section{Системы управления версиями}

\subsection{Назначение систем управления версиями}

\textbf{Система управления версиями} --- это программная система, предназначенная для хранения истории изменений файлов, координации совместной работы и восстановления предыдущих состояний проекта.

Системы управления версиями позволяют:

\begin{itemize}
    \item хранить историю изменений;
    \item сравнивать версии файлов;
    \item возвращаться к предыдущим состояниям;
    \item параллельно работать над разными задачами;
    \item объединять изменения нескольких разработчиков;
    \item фиксировать авторство изменений;
    \item создавать ветки разработки;
    \item помечать релизы тегами.
\end{itemize}

\subsection{Основные понятия}

\textbf{Репозиторий} --- хранилище проекта и его истории.

\textbf{Рабочая копия} --- локальный набор файлов, с которым непосредственно работает разработчик.

\textbf{Коммит} --- зафиксированное изменение состояния проекта.

\textbf{Ветка} --- именованный указатель на некоторую линию разработки.

\textbf{Слияние} --- операция объединения изменений из разных веток.

\textbf{Конфликт} --- ситуация, когда система не может автоматически объединить изменения.

\textbf{Тег} --- именованная метка на конкретный коммит, обычно используемая для обозначения релизов.

\subsection{Локальные, централизованные и распределённые системы}

Исторически можно выделить три типа систем управления версиями:

\begin{enumerate}
    \item локальные;
    \item централизованные;
    \item распределённые.
\end{enumerate}

\subsubsection{Локальные системы управления версиями}

В локальной системе история изменений хранится на одной машине.

Простейший пример неформального локального версионирования:

\begin{lstlisting}
project_final.zip
project_final_2.zip
project_final_real.zip
project_final_real_fixed.zip
\end{lstlisting}

Недостатки такого подхода:

\begin{itemize}
    \item сложно понять, какая версия актуальна;
    \item нет полноценной истории изменений;
    \item высок риск потери данных;
    \item неудобна совместная работа.
\end{itemize}

\subsubsection{Централизованные системы управления версиями}

\textbf{Централизованная система управления версиями} хранит основной репозиторий на центральном сервере. Разработчики получают рабочие копии и отправляют изменения обратно на сервер.

Примеры:

\begin{itemize}
    \item CVS;
    \item Subversion;
    \item Perforce;
    \item Team Foundation Version Control.
\end{itemize}

Схема централизованной модели:

\begin{center}
\begin{tikzpicture}[node distance=2.5cm]
    \node[draw, rectangle, rounded corners] (server) {Центральный репозиторий};
    \node[draw, rectangle, rounded corners, below left of=server, xshift=-1.5cm] (dev1) {Разработчик 1};
    \node[draw, rectangle, rounded corners, below of=server] (dev2) {Разработчик 2};
    \node[draw, rectangle, rounded corners, below right of=server, xshift=1.5cm] (dev3) {Разработчик 3};

    \draw[<->] (server) -- (dev1);
    \draw[<->] (server) -- (dev2);
    \draw[<->] (server) -- (dev3);
\end{tikzpicture}
\end{center}

Преимущества централизованной модели:

\begin{itemize}
    \item простая модель доступа;
    \item легко администрировать права;
    \item существует единый источник истины;
    \item удобно контролировать состояние проекта.
\end{itemize}

Недостатки:

\begin{itemize}
    \item зависимость от центрального сервера;
    \item при недоступности сервера затруднена работа;
    \item сервер является единой точкой отказа;
    \item многие операции требуют сетевого соединения;
    \item ветвление и слияние часто тяжелее, чем в распределённых системах.
\end{itemize}

\subsubsection{Распределённые системы управления версиями}

\textbf{Распределённая система управления версиями} хранит полную копию репозитория и истории у каждого разработчика.

Примеры:

\begin{itemize}
    \item Git;
    \item Mercurial;
    \item Bazaar;
    \item Fossil.
\end{itemize}

Схема распределённой модели:

\begin{center}
\begin{tikzpicture}[node distance=3cm]
    \node[draw, rectangle, rounded corners] (remote) {Удалённый репозиторий};
    \node[draw, rectangle, rounded corners, below left of=remote, xshift=-2cm] (dev1) {Полный репозиторий 1};
    \node[draw, rectangle, rounded corners, below of=remote] (dev2) {Полный репозиторий 2};
    \node[draw, rectangle, rounded corners, below right of=remote, xshift=2cm] (dev3) {Полный репозиторий 3};

    \draw[<->] (remote) -- (dev1);
    \draw[<->] (remote) -- (dev2);
    \draw[<->] (remote) -- (dev3);
    \draw[<->, dashed] (dev1) -- (dev2);
    \draw[<->, dashed] (dev2) -- (dev3);
\end{tikzpicture}
\end{center}

Преимущества распределённой модели:

\begin{itemize}
    \item большинство операций выполняется локально;
    \item можно работать без сети;
    \item каждый клон является резервной копией проекта;
    \item удобно создавать ветки;
    \item проще организовать разные модели совместной работы;
    \item выше гибкость при интеграции изменений.
\end{itemize}

Недостатки:

\begin{itemize}
    \item сложнее концептуальная модель;
    \item необходимо понимать различие между локальным и удалённым состоянием;
    \item требуется дисциплина при работе с ветками;
    \item возможны сложные истории слияний.
\end{itemize}

\subsection{Сравнение централизованного и распределённого управления версиями}

\begin{center}
\begin{tabular}{|p{4cm}|p{5cm}|p{5cm}|}
\hline
\textbf{Критерий} & \textbf{Централизованная модель} & \textbf{Распределённая модель} \\
\hline
Хранение истории & На центральном сервере & У каждого разработчика \\
\hline
Работа без сети & Ограничена & Возможна для большинства операций \\
\hline
Единая точка отказа & Центральный сервер & Обычно отсутствует \\
\hline
Скорость операций & Часто зависит от сети & Большинство операций локальны \\
\hline
Ветвление & Может быть дорогим или неудобным & Обычно лёгкое и быстрое \\
\hline
Администрирование & Проще централизованно контролировать & Гибче, но сложнее модель доступа \\
\hline
Примеры & SVN, CVS, Perforce & Git, Mercurial \\
\hline
\end{tabular}
\end{center}

\section{Git}

\subsection{Общая характеристика Git}

\textbf{Git} --- распределённая система управления версиями, ориентированная на скорость, надёжность, поддержку ветвления и распределённой совместной работы.

Git хранит не просто последовательность различий между файлами, а набор снимков состояния проекта. Если файл не изменился, Git может ссылаться на уже существующее содержимое.

\subsection{Основные области Git}

В Git удобно выделять несколько областей:

\begin{enumerate}
    \item \textbf{Рабочая директория} --- файлы, с которыми работает пользователь.
    \item \textbf{Индекс} --- промежуточная область подготовки коммита.
    \item \textbf{Локальный репозиторий} --- база объектов и история коммитов.
    \item \textbf{Удалённый репозиторий} --- репозиторий на сервере или другой машине.
\end{enumerate}

Схема движения изменений:

\begin{center}
\begin{tikzpicture}[node distance=3cm]
    \node[draw, rectangle, rounded corners] (work) {Рабочая директория};
    \node[draw, rectangle, rounded corners, right of=work] (index) {Индекс};
    \node[draw, rectangle, rounded corners, right of=index] (local) {Локальный репозиторий};
    \node[draw, rectangle, rounded corners, right of=local] (remote) {Удалённый репозиторий};

    \draw[->] (work) -- node[above]{\texttt{git add}} (index);
    \draw[->] (index) -- node[above]{\texttt{git commit}} (local);
    \draw[->] (local) -- node[above]{\texttt{git push}} (remote);
    \draw[->, bend right=25] (remote) to node[below]{\texttt{git pull/fetch}} (local);
\end{tikzpicture}
\end{center}

\subsection{Базовый цикл работы в Git}

Алгоритм обычной работы с Git:

\begin{enumerate}
    \item Получить актуальное состояние проекта:
\begin{lstlisting}
git pull
\end{lstlisting}

    \item Изменить файлы в рабочей директории.

    \item Посмотреть состояние:
\begin{lstlisting}
git status
\end{lstlisting}

    \item Добавить изменения в индекс:
\begin{lstlisting}
git add file.py
\end{lstlisting}

    \item Создать коммит:
\begin{lstlisting}
git commit -m "Add input validation"
\end{lstlisting}

    \item Отправить изменения в удалённый репозиторий:
\begin{lstlisting}
git push
\end{lstlisting}
\end{enumerate}

\subsection{Коммит}

\textbf{Коммит} в Git содержит:

\begin{itemize}
    \item ссылку на снимок дерева файлов;
    \item ссылку на одного или нескольких родителей;
    \item автора;
    \item дату;
    \item сообщение;
    \item хеш-идентификатор.
\end{itemize}

Если коммит имеет одного родителя, это обычный коммит. Если несколько родителей --- это коммит слияния.

\subsection{Ветка}

\textbf{Ветка} в Git --- это подвижный указатель на коммит.

Пример:

\begin{lstlisting}
git branch feature-login
git switch feature-login
\end{lstlisting}

или короче:

\begin{lstlisting}
git switch -c feature-login
\end{lstlisting}

Схема:

\begin{center}
\begin{tikzpicture}
    \node[circle, draw] (A) at (0,0) {A};
    \node[circle, draw] (B) at (2,0) {B};
    \node[circle, draw] (C) at (4,0) {C};
    \node[circle, draw] (D) at (6,1.2) {D};
    \node[circle, draw] (E) at (6,-1.2) {E};

    \draw[->] (A) -- (B);
    \draw[->] (B) -- (C);
    \draw[->] (C) -- (D);
    \draw[->] (C) -- (E);

    \node[right of=D, xshift=1cm] {feature};
    \node[right of=E, xshift=1cm] {main};
\end{tikzpicture}
\end{center}

\subsection{Слияние веток}

\textbf{Слияние} --- операция объединения изменений из одной ветки в другую.

Пример:

\begin{lstlisting}
git switch main
git merge feature-login
\end{lstlisting}

\subsubsection{Fast-forward merge}

Если целевая ветка не содержит новых коммитов после момента расхождения, Git может просто передвинуть указатель ветки вперёд.

До слияния:

\[
A \rightarrow B \rightarrow C \quad (\texttt{main})
\]
\[
C \rightarrow D \rightarrow E \quad (\texttt{feature})
\]

После fast-forward:

\[
A \rightarrow B \rightarrow C \rightarrow D \rightarrow E \quad (\texttt{main})
\]

\subsubsection{Трёхстороннее слияние}

Если обе ветки изменялись независимо, Git использует:

\begin{itemize}
    \item общего предка;
    \item текущую версию;
    \item сливаемую версию.
\end{itemize}

Алгоритм трёхстороннего слияния:

\begin{enumerate}
    \item Найти общего предка двух веток.
    \item Сравнить изменения от предка до первой ветки.
    \item Сравнить изменения от предка до второй ветки.
    \item Объединить независимые изменения автоматически.
    \item Если изменения затрагивают одни и те же строки несовместимым образом, зафиксировать конфликт.
    \item После разрешения конфликтов создать коммит слияния.
\end{enumerate}

Схема:

\begin{center}
\begin{tikzpicture}
    \node[circle, draw] (A) at (0,0) {A};
    \node[circle, draw] (B) at (2,0) {B};
    \node[circle, draw] (C) at (4,1.2) {C};
    \node[circle, draw] (D) at (4,-1.2) {D};
    \node[circle, draw] (M) at (6,0) {M};

    \draw[->] (A) -- (B);
    \draw[->] (B) -- (C);
    \draw[->] (B) -- (D);
    \draw[->] (C) -- (M);
    \draw[->] (D) -- (M);

    \node[below of=B, yshift=1cm] {общий предок};
    \node[right of=M, xshift=0.8cm] {merge commit};
\end{tikzpicture}
\end{center}

\subsection{Конфликты слияния}

Конфликт возникает, когда Git не может автоматически определить итоговую версию файла.

Пример конфликта:

\begin{lstlisting}
<<<<<<< HEAD
print("Hello from main")
=======
print("Hello from feature")
>>>>>>> feature
\end{lstlisting}

Алгоритм разрешения конфликта:

\begin{enumerate}
    \item Открыть конфликтующий файл.
    \item Найти конфликтные блоки.
    \item Выбрать одну из версий или вручную объединить изменения.
    \item Удалить служебные маркеры:
\begin{lstlisting}
<<<<<<<
=======
>>>>>>>
\end{lstlisting}
    \item Добавить исправленный файл в индекс:
\begin{lstlisting}
git add file.py
\end{lstlisting}
    \item Завершить слияние:
\begin{lstlisting}
git commit
\end{lstlisting}
\end{enumerate}

\subsection{Rebase}

\textbf{Rebase} --- операция переноса коммитов одной ветки на новое основание.

Пример:

\begin{lstlisting}
git switch feature
git rebase main
\end{lstlisting}

До rebase:

\[
A \rightarrow B \rightarrow C \quad (\texttt{main})
\]
\[
B \rightarrow D \rightarrow E \quad (\texttt{feature})
\]

После rebase:

\[
A \rightarrow B \rightarrow C \rightarrow D' \rightarrow E' \quad (\texttt{feature})
\]

Коммиты $D'$ и $E'$ являются новыми коммитами с теми же изменениями, но с другими родителями и другими хешами.

Основное правило:

\begin{quote}
Не следует выполнять rebase опубликованных коммитов, если с ними уже работают другие разработчики.
\end{quote}

\subsection{Теги}

\textbf{Тег} --- ссылка на конкретный коммит. Обычно теги используются для обозначения релизов.

Пример:

\begin{lstlisting}
git tag v1.0.0
git push origin v1.0.0
\end{lstlisting}

Аннотированный тег:

\begin{lstlisting}
git tag -a v1.0.0 -m "Release version 1.0.0"
\end{lstlisting}

\section{Git Workflow}

\subsection{Понятие Git Workflow}

\textbf{Git Workflow} --- это соглашение о том, как команда использует Git: какие ветки создаются, как именуются, как проходят ревью, как выполняется интеграция и выпуск релизов.

Git сам по себе не навязывает единственный процесс. Поэтому разные команды применяют разные модели:

\begin{itemize}
    \item централизованный workflow;
    \item feature branch workflow;
    \item Git Flow;
    \item GitHub Flow;
    \item GitLab Flow;
    \item trunk-based development.
\end{itemize}

\subsection{Централизованный workflow в Git}

Несмотря на распределённую природу Git, его можно использовать в централизованном стиле.

Основная идея:

\begin{itemize}
    \item есть одна главная ветка, например \texttt{main};
    \item все разработчики клонируют один удалённый репозиторий;
    \item изменения отправляются в центральный репозиторий;
    \item перед отправкой нужно получить актуальное состояние.
\end{itemize}

Алгоритм:

\begin{enumerate}
    \item Клонировать репозиторий:
\begin{lstlisting}
git clone https://example.com/project.git
\end{lstlisting}

    \item Получить изменения:
\begin{lstlisting}
git pull origin main
\end{lstlisting}

    \item Внести изменения.

    \item Зафиксировать изменения:
\begin{lstlisting}
git add .
git commit -m "Fix typo"
\end{lstlisting}

    \item Отправить изменения:
\begin{lstlisting}
git push origin main
\end{lstlisting}
\end{enumerate}

Преимущество --- простота. Недостаток --- высокая вероятность конфликтов и отсутствие изоляции задач.

\subsection{Feature Branch Workflow}

\textbf{Feature Branch Workflow} --- модель, при которой каждая задача разрабатывается в отдельной ветке.

Алгоритм:

\begin{enumerate}
    \item Обновить основную ветку:
\begin{lstlisting}
git switch main
git pull
\end{lstlisting}

    \item Создать ветку задачи:
\begin{lstlisting}
git switch -c feature/payment
\end{lstlisting}

    \item Выполнить разработку и создать коммиты:
\begin{lstlisting}
git add .
git commit -m "Add payment form"
\end{lstlisting}

    \item Отправить ветку:
\begin{lstlisting}
git push -u origin feature/payment
\end{lstlisting}

    \item Создать запрос на слияние.
    \item Провести ревью.
    \item Выполнить автоматические проверки.
    \item Слить ветку в \texttt{main}.
    \item Удалить ветку.
\end{enumerate}

Преимущества:

\begin{itemize}
    \item изоляция задач;
    \item удобное ревью;
    \item можно запускать CI для каждой ветки;
    \item история изменений становится понятнее.
\end{itemize}

\subsection{Git Flow}

\textbf{Git Flow} --- модель ветвления, в которой используются несколько типов долгоживущих и короткоживущих веток.

Основные ветки:

\begin{itemize}
    \item \texttt{main} --- стабильные релизы;
    \item \texttt{develop} --- основная ветка разработки.
\end{itemize}

Вспомогательные ветки:

\begin{itemize}
    \item \texttt{feature/*} --- разработка новых возможностей;
    \item \texttt{release/*} --- подготовка релиза;
    \item \texttt{hotfix/*} --- срочные исправления в промышленной версии.
\end{itemize}

Схема Git Flow:

\begin{center}
\begin{tikzpicture}
    \node[circle, draw] (M1) at (0,0) {};
    \node[circle, draw] (M2) at (3,0) {};
    \node[circle, draw] (M3) at (6,0) {};
    \node[right] at (6.5,0) {\texttt{main}};

    \node[circle, draw] (D1) at (0,-1.5) {};
    \node[circle, draw] (D2) at (2,-1.5) {};
    \node[circle, draw] (D3) at (4,-1.5) {};
    \node[circle, draw] (D4) at (6,-1.5) {};
    \node[right] at (6.5,-1.5) {\texttt{develop}};

    \node[circle, draw] (F1) at (2,-3) {};
    \node[circle, draw] (F2) at (4,-3) {};
    \node[right] at (4.5,-3) {\texttt{feature}};

    \draw[->] (M1) -- (M2);
    \draw[->] (M2) -- (M3);

    \draw[->] (D1) -- (D2);
    \draw[->] (D2) -- (D3);
    \draw[->] (D3) -- (D4);

    \draw[->] (D2) -- (F1);
    \draw[->] (F1) -- (F2);
    \draw[->] (F2) -- (D3);

    \draw[->] (D4) -- (M3);
\end{tikzpicture}
\end{center}

Алгоритм разработки новой возможности:

\begin{enumerate}
    \item Создать ветку от \texttt{develop}:
\begin{lstlisting}
git switch develop
git pull
git switch -c feature/search
\end{lstlisting}

    \item Реализовать функциональность.
    \item Слить ветку обратно в \texttt{develop}:
\begin{lstlisting}
git switch develop
git merge feature/search
\end{lstlisting}
\end{enumerate}

Алгоритм выпуска релиза:

\begin{enumerate}
    \item Создать ветку релиза:
\begin{lstlisting}
git switch develop
git switch -c release/1.2.0
\end{lstlisting}

    \item Исправлять только ошибки и релизные проблемы.
    \item Слить релиз в \texttt{main}.
    \item Поставить тег:
\begin{lstlisting}
git tag -a v1.2.0 -m "Release 1.2.0"
\end{lstlisting}

    \item Слить релизную ветку обратно в \texttt{develop}.
\end{enumerate}

Алгоритм hotfix:

\begin{enumerate}
    \item Создать ветку от \texttt{main}:
\begin{lstlisting}
git switch main
git switch -c hotfix/critical-bug
\end{lstlisting}

    \item Исправить ошибку.
    \item Слить исправление в \texttt{main}.
    \item Поставить новый тег.
    \item Слить исправление в \texttt{develop}.
\end{enumerate}

Git Flow хорошо подходит для проектов с явно выраженными релизными циклами. Для проектов с частой поставкой он может быть избыточным.

\section{GitHub Workflow}

\subsection{GitHub как платформа совместной разработки}

GitHub --- это платформа для размещения Git-репозиториев и организации совместной разработки.

GitHub предоставляет:

\begin{itemize}
    \item удалённые Git-репозитории;
    \item pull request;
    \item code review;
    \item issues;
    \item project boards;
    \item GitHub Actions;
    \item protected branches;
    \item releases;
    \item packages;
    \item wiki и документацию.
\end{itemize}

\subsection{GitHub Flow}

\textbf{GitHub Flow} --- простой workflow, основанный на короткоживущих ветках, pull request и непрерывной интеграции.

Основная идея:

\begin{itemize}
    \item ветка \texttt{main} всегда должна быть готова к развёртыванию;
    \item любая новая работа выполняется в отдельной ветке;
    \item изменения обсуждаются через pull request;
    \item перед слиянием запускаются проверки;
    \item после слияния изменения могут быть развёрнуты.
\end{itemize}

\subsection{Алгоритм GitHub Flow}

\begin{enumerate}
    \item Создать ветку от \texttt{main}:
\begin{lstlisting}
git switch main
git pull
git switch -c feature/user-profile
\end{lstlisting}

    \item Выполнить изменения и создать коммиты:
\begin{lstlisting}
git add .
git commit -m "Add user profile page"
\end{lstlisting}

    \item Отправить ветку на GitHub:
\begin{lstlisting}
git push -u origin feature/user-profile
\end{lstlisting}

    \item Создать pull request.

    \item Обсудить изменения и выполнить code review.

    \item Дождаться успешного прохождения CI-проверок.

    \item При необходимости внести исправления.

    \item Слить pull request в \texttt{main}.

    \item Удалить feature-ветку.

    \item Развернуть изменения.
\end{enumerate}

\subsection{Pull request}

\textbf{Pull request} --- запрос на включение изменений из одной ветки в другую, сопровождаемый обсуждением, проверками и ревью.

Pull request обычно содержит:

\begin{itemize}
    \item описание задачи;
    \item список изменённых файлов;
    \item набор коммитов;
    \item комментарии участников;
    \item результаты автоматических проверок;
    \item решение о принятии или отклонении изменений.
\end{itemize}

\subsection{Code review}

\textbf{Code review} --- проверка изменений другими разработчиками.

Цели code review:

\begin{itemize}
    \item обнаружить ошибки;
    \item проверить читаемость и сопровождаемость кода;
    \item обеспечить соответствие архитектуре;
    \item распространить знания внутри команды;
    \item обеспечить единый стиль кода;
    \item предотвратить попадание небезопасного кода в основную ветку.
\end{itemize}

\subsection{Protected branches}

\textbf{Protected branch} --- защищённая ветка, для которой задаются дополнительные правила.

Например, для \texttt{main} можно потребовать:

\begin{itemize}
    \item запрет прямого \texttt{push};
    \item обязательный pull request;
    \item минимум одно одобрение ревьюера;
    \item успешное прохождение CI;
    \item актуальность ветки перед слиянием;
    \item подписанные коммиты.
\end{itemize}

\section{CI/CD}

\subsection{Общее понятие}

\textbf{CI/CD} --- совокупность практик автоматизации интеграции, проверки, сборки и поставки программного обеспечения.

Аббревиатура обычно расшифровывается как:

\begin{itemize}
    \item \textbf{CI} --- Continuous Integration, непрерывная интеграция;
    \item \textbf{CD} --- Continuous Delivery, непрерывная доставка;
    \item \textbf{CD} --- Continuous Deployment, непрерывное развёртывание.
\end{itemize}

\subsection{Continuous Integration}

\textbf{Непрерывная интеграция} --- практика частого объединения изменений разработчиков в общую ветку с автоматической сборкой и тестированием.

Основные идеи CI:

\begin{itemize}
    \item изменения интегрируются часто;
    \item сборка выполняется автоматически;
    \item тесты запускаются автоматически;
    \item ошибки обнаруживаются как можно раньше;
    \item основная ветка поддерживается в рабочем состоянии.
\end{itemize}

Типичный CI-процесс:

\begin{enumerate}
    \item Разработчик отправляет коммит или создаёт pull request.
    \item CI-сервер получает событие.
    \item Создаётся чистое окружение.
    \item Загружается исходный код.
    \item Устанавливаются зависимости.
    \item Выполняется статический анализ.
    \item Запускается сборка.
    \item Запускаются тесты.
    \item Формируется отчёт.
    \item Результат передаётся в систему управления версиями.
\end{enumerate}

\subsection{Continuous Delivery}

\textbf{Непрерывная доставка} --- практика, при которой каждое изменение, прошедшее автоматические проверки, может быть развёрнуто в эксплуатацию, но само развёртывание обычно запускается вручную.

Иными словами:

\[
\text{Continuous Delivery} = \text{CI} + \text{автоматизированная подготовка релиза}
\]

При continuous delivery система всегда находится в состоянии, пригодном для выпуска.

\subsection{Continuous Deployment}

\textbf{Непрерывное развёртывание} --- практика, при которой каждое изменение, прошедшее все проверки, автоматически разворачивается в целевую среду.

Иными словами:

\[
\text{Continuous Deployment} = \text{Continuous Delivery} + \text{автоматический production deploy}
\]

Различие:

\begin{center}
\begin{tabular}{|p{5cm}|p{8cm}|}
\hline
\textbf{Практика} & \textbf{Смысл} \\
\hline
Continuous Integration & Автоматически собрать и проверить изменения \\
\hline
Continuous Delivery & Автоматически подготовить поставку, но выпуск может быть ручным \\
\hline
Continuous Deployment & Автоматически развернуть изменения после успешных проверок \\
\hline
\end{tabular}
\end{center}

\subsection{Pipeline}

\textbf{Pipeline} --- автоматизированная последовательность стадий, через которые проходит изменение от коммита до поставки.

Пример pipeline:

\[
\text{commit} \rightarrow \text{lint} \rightarrow \text{test} \rightarrow \text{build} \rightarrow \text{package} \rightarrow \text{deploy}
\]

Схема:

\begin{center}
\begin{tikzpicture}[node distance=2.5cm]
    \node[draw, rectangle, rounded corners] (commit) {Commit};
    \node[draw, rectangle, rounded corners, right of=commit] (lint) {Lint};
    \node[draw, rectangle, rounded corners, right of=lint] (test) {Test};
    \node[draw, rectangle, rounded corners, right of=test] (build) {Build};
    \node[draw, rectangle, rounded corners, below of=build] (deploy) {Deploy};

    \draw[->] (commit) -- (lint);
    \draw[->] (lint) -- (test);
    \draw[->] (test) -- (build);
    \draw[->] (build) -- (deploy);
\end{tikzpicture}
\end{center}

\subsection{Основные стадии CI/CD}

\subsubsection{Checkout}

Получение исходного кода из репозитория.

\begin{lstlisting}
git checkout main
\end{lstlisting}

\subsubsection{Install}

Установка зависимостей.

Например, для Python:

\begin{lstlisting}
pip install -r requirements.txt
\end{lstlisting}

Для Node.js:

\begin{lstlisting}
npm ci
\end{lstlisting}

\subsubsection{Lint}

Проверка стиля и потенциальных ошибок.

\begin{lstlisting}
flake8 src tests
\end{lstlisting}

\subsubsection{Test}

Запуск автоматических тестов.

\begin{lstlisting}
pytest
\end{lstlisting}

\subsubsection{Build}

Сборка приложения или артефакта.

\begin{lstlisting}
docker build -t app:latest .
\end{lstlisting}

\subsubsection{Package}

Упаковка результата.

Примеры артефактов:

\begin{itemize}
    \item Docker-образ;
    \item JAR-файл;
    \item wheel-пакет;
    \item архив;
    \item бинарный исполняемый файл.
\end{itemize}

\subsubsection{Deploy}

Развёртывание в среду:

\begin{itemize}
    \item development;
    \item testing;
    \item staging;
    \item production.
\end{itemize}

\subsection{Минимальный пример GitHub Actions}

GitHub Actions позволяет описывать workflows в YAML-файлах, расположенных в каталоге:

\begin{lstlisting}
.github/workflows/
\end{lstlisting}

Минимальный пример CI для Python-проекта:

\begin{lstlisting}
name: CI

on:
  pull_request:
  push:
    branches:
      - main

jobs:
  test:
    runs-on: ubuntu-latest

    steps:
      - name: Checkout repository
        uses: actions/checkout@v4

      - name: Set up Python
        uses: actions/setup-python@v5
        with:
          python-version: "3.12"

      - name: Install dependencies
        run: pip install -r requirements.txt

      - name: Run tests
        run: pytest
\end{lstlisting}

Пояснение:

\begin{itemize}
    \item \texttt{name} задаёт имя workflow;
    \item \texttt{on} описывает события запуска;
    \item \texttt{jobs} содержит задания;
    \item \texttt{runs-on} задаёт окружение выполнения;
    \item \texttt{steps} описывает последовательность шагов;
    \item \texttt{uses} подключает готовое действие;
    \item \texttt{run} выполняет shell-команду.
\end{itemize}

\subsection{Минимальный пример CI/CD pipeline}

Пример workflow, который тестирует приложение и затем собирает Docker-образ:

\begin{lstlisting}
name: CI-CD

on:
  push:
    branches:
      - main

jobs:
  test:
    runs-on: ubuntu-latest

    steps:
      - uses: actions/checkout@v4

      - name: Run tests
        run: |
          pip install -r requirements.txt
          pytest

  build:
    needs: test
    runs-on: ubuntu-latest

    steps:
      - uses: actions/checkout@v4

      - name: Build Docker image
        run: docker build -t my-app:latest .
\end{lstlisting}

Здесь:

\begin{itemize}
    \item job \texttt{test} запускает тесты;
    \item job \texttt{build} зависит от \texttt{test};
    \item если тесты падают, сборка не выполняется.
\end{itemize}

\subsection{Алгоритм работы CI/CD-системы}

\begin{enumerate}
    \item Разработчик создаёт изменение в репозитории.
    \item Система управления версиями генерирует событие.
    \item CI/CD-система получает событие через webhook или внутренний механизм.
    \item Запускается pipeline.
    \item Pipeline создаёт изолированное окружение.
    \item Исходный код загружается из репозитория.
    \item Выполняются команды сборки и проверки.
    \item При успехе создаются артефакты.
    \item Артефакты сохраняются в хранилище.
    \item При необходимости выполняется развёртывание.
    \item Результат возвращается в pull request или commit status.
\end{enumerate}

\subsection{Популярные CI/CD-системы}

Примеры систем:

\begin{itemize}
    \item GitHub Actions;
    \item GitLab CI/CD;
    \item Jenkins;
    \item TeamCity;
    \item CircleCI;
    \item Travis CI;
    \item Azure Pipelines;
    \item Bitbucket Pipelines.
\end{itemize}

\subsection{Артефакты}

\textbf{Артефакт} --- результат выполнения pipeline, который можно сохранить, передать или развернуть.

Примеры:

\begin{itemize}
    \item исполняемый файл;
    \item Docker-образ;
    \item отчёт о тестах;
    \item отчёт о покрытии;
    \item архив собранного приложения;
    \item скомпилированная документация.
\end{itemize}

\subsection{Окружения}

Обычно выделяют несколько окружений:

\begin{itemize}
    \item \textbf{local} --- машина разработчика;
    \item \textbf{development} --- среда разработки;
    \item \textbf{testing} --- среда тестирования;
    \item \textbf{staging} --- предпромышленная среда;
    \item \textbf{production} --- промышленная среда.
\end{itemize}

Хорошая практика состоит в том, чтобы окружения были максимально похожи друг на друга.

\subsection{Стратегии развёртывания}

\subsubsection{Rolling deployment}

\textbf{Rolling deployment} --- постепенная замена старых экземпляров приложения новыми.

Преимущество: нет полной остановки сервиса.

Недостаток: некоторое время одновременно существуют разные версии приложения.

\subsubsection{Blue-green deployment}

\textbf{Blue-green deployment} использует две среды:

\begin{itemize}
    \item blue --- текущая промышленная версия;
    \item green --- новая версия.
\end{itemize}

Алгоритм:

\begin{enumerate}
    \item Развернуть новую версию в неактивной среде.
    \item Проверить её.
    \item Переключить трафик на новую среду.
    \item При ошибке быстро вернуть трафик обратно.
\end{enumerate}

\subsubsection{Canary deployment}

\textbf{Canary deployment} --- постепенное направление небольшой части пользователей на новую версию.

Алгоритм:

\begin{enumerate}
    \item Развернуть новую версию для малой доли пользователей.
    \item Собрать метрики и ошибки.
    \item Если всё хорошо, увеличить долю трафика.
    \item При проблемах откатить новую версию.
\end{enumerate}

\subsection{Преимущества CI/CD}

CI/CD даёт следующие преимущества:

\begin{itemize}
    \item раннее обнаружение ошибок;
    \item сокращение ручного труда;
    \item воспроизводимость сборок;
    \item ускорение поставки изменений;
    \item повышение качества кода;
    \item прозрачность процесса разработки;
    \item уменьшение риска релизов;
    \item возможность частых малых поставок.
\end{itemize}

\subsection{Типичные ошибки при внедрении CI/CD}

\begin{itemize}
    \item отсутствие автоматических тестов;
    \item нестабильные тесты;
    \item слишком долгий pipeline;
    \item хранение секретов в репозитории;
    \item отсутствие изолированных окружений;
    \item ручные неповторяемые шаги;
    \item отсутствие rollback-стратегии;
    \item игнорирование failed build;
    \item слишком долгоживущие ветки.
\end{itemize}

\section{Связь Git Workflow, GitHub Workflow и CI/CD}

Современный процесс разработки обычно объединяет все рассмотренные элементы.

Типичный цикл:

\begin{enumerate}
    \item Задача регистрируется в issue tracker.
    \item Разработчик создаёт feature-ветку.
    \item Изменения фиксируются коммитами.
    \item Ветка отправляется в удалённый репозиторий.
    \item Создаётся pull request.
    \item Запускается CI.
    \item Выполняется code review.
    \item После успешных проверок изменения сливаются в \texttt{main}.
    \item Запускается pipeline поставки.
    \item Создаётся артефакт.
    \item Артефакт разворачивается в нужную среду.
    \item Релиз помечается тегом.
\end{enumerate}

Схема:

\begin{center}
\begin{tikzpicture}[node distance=2.8cm]
    \node[draw, rectangle, rounded corners] (issue) {Issue};
    \node[draw, rectangle, rounded corners, right of=issue] (branch) {Branch};
    \node[draw, rectangle, rounded corners, right of=branch] (pr) {Pull Request};
    \node[draw, rectangle, rounded corners, below of=pr] (ci) {CI};
    \node[draw, rectangle, rounded corners, left of=ci] (review) {Review};
    \node[draw, rectangle, rounded corners, left of=review] (merge) {Merge};
    \node[draw, rectangle, rounded corners, below of=merge] (cd) {CD};
    \node[draw, rectangle, rounded corners, right of=cd] (release) {Release};

    \draw[->] (issue) -- (branch);
    \draw[->] (branch) -- (pr);
    \draw[->] (pr) -- (ci);
    \draw[->] (ci) -- (review);
    \draw[->] (review) -- (merge);
    \draw[->] (merge) -- (cd);
    \draw[->] (cd) -- (release);
\end{tikzpicture}
\end{center}

\section{Краткие выводы}

Конфигурационное управление обеспечивает контроль состава и состояния программной системы. Системы управления версиями являются одним из ключевых инструментов конфигурационного управления. Централизованные системы используют единый центральный репозиторий, а распределённые предоставляют каждому разработчику полную копию истории.

Git является распределённой системой управления версиями, в которой ветки и слияния являются дешёвыми и удобными операциями. Git Workflow описывает правила работы команды с ветками, коммитами, ревью и релизами. GitHub Workflow дополняет Git механизмами pull request, code review, protected branches и автоматических проверок.

CI/CD автоматизирует сборку, тестирование и поставку программного обеспечения. В связке с Git и GitHub CI/CD позволяет строить воспроизводимый, контролируемый и быстрый процесс разработки, в котором каждое изменение проходит проверку перед попаданием в основную ветку и релиз.

\end{document}
